\documentclass[10pt,letterpaper]{article}

\usepackage[T1]{fontenc}
\usepackage[utf8]{inputenc}
\usepackage{newtxtext}
\usepackage[margin=0.78in]{geometry}
\usepackage{array,booktabs,tabularx}
\usepackage{xcolor}
\usepackage{enumitem}
\usepackage{titlesec}
\usepackage{microtype}
\usepackage[hidelinks]{hyperref}
\usepackage[skip=0.42em]{parskip}

\definecolor{reviewblue}{RGB}{30,70,110}
\newcolumntype{Y}{>{\raggedright\arraybackslash}X}
\setlist[itemize]{leftmargin=1.45em,itemsep=0.2em,topsep=0.2em}
\setlength{\parindent}{0pt}
\titlespacing*{\section}{0pt}{0.9em}{0.35em}

\hypersetup{
  pdftitle={Response to Reviewer Comments for ACOMPA 2026},
  pdfauthor={Nhat Phat Huynh; Dong Quan Nguyen; Quan Le; Thanh Hung Le; Minh Son Nguyen},
  pdfsubject={Author response for A New IoT Platform Based on PQC Algorithms},
  pdfkeywords={ACOMPA 2026, author response, post-quantum cryptography, Internet of Things}
}

\begin{document}

\begin{center}
{\Large\bfseries Response to Reviewer Comments}\par
\vspace{0.5em}
{\bfseries ACOMPA 2026}\par
\vspace{0.4em}
\fbox{\bfseries Working draft---final figure artwork pending}\par
\vspace{0.4em}
{\large\itshape A New IoT Platform Based on PQC Algorithms:\ A Solution for Quantum Computing Threats}\par
\vspace{0.5em}
Nhat Phat Huynh, Dong Quan Nguyen, Quan Le, Thanh Hung Le, and Minh Son Nguyen
\end{center}

\vspace{0.75em}

We thank the reviewer for the careful assessment and for recognizing the practical, architecture-driven treatment of post-quantum cryptography (PQC) migration in heterogeneous Internet of Things (IoT) systems. We agree that the figures in the submitted manuscript did not meet publication quality. This response is a working draft because the two figure files currently included in the revised manuscript are placeholders extracted from the latest DOCX version. The figure-quality response must be updated after the authors supply the final redrawn artwork. The remaining responses identify the revisions already present in the manuscript and their locations.

\section*{Response to the overall assessment}

\textcolor{reviewblue}{\textbf{Reviewer assessment:}} \emph{Borderline Accept. The paper's core strength lies in its practical, architecture-driven approach to migrating heterogeneous IoT ecosystems to PQC. The study introduces a strategic algorithm-placement model across device tiers.}

\textbf{Response.} We appreciate this assessment. The revised manuscript retains the architecture-driven scope and makes the placement argument easier to verify. Section III-A now presents the node, gateway, and server responsibilities as an explicit tiered placement model, and Table I states the placement and design rationale for ML-KEM, ML-DSA, SLH-DSA, and hybrid mode. The Introduction states the contribution in the same terms, while the Conclusion limits the claim to a migration-oriented architecture supported by the reported prototype measurements.

\section*{Response to the figure-quality comment}

\textcolor{reviewblue}{\textbf{Reviewer comment:}} \emph{I am mainly concerned about the AI-generated figures (Fig. 1 and Fig. 4). The human-generated figures (Fig. 2 and Fig. 3) are of low quality and therefore must be improved.}

\textbf{Response.} We agree. The four-figure structure of the submitted manuscript has been consolidated into two figures to remove overlap between the overview, hardware, sequence, and deployment views. At the authors' instruction, however, the current compiled manuscript uses the two raster images embedded in DOCX V4.0 as temporary placeholders. This working draft therefore does not claim that the reviewer's figure-quality concern has been fully resolved. Before camera-ready submission, both placeholders must be replaced by the final author-supplied redraws, and this response must be updated to state the actual production method and file format.

The consolidation was deliberate:

\begin{itemize}
  \item \textbf{Fig. 1, ``End-node architecture for PQC-enabled IoT devices,''} carries the end-node processing, secure-boot, memory, storage, network, and peripheral content discussed in Section III-C. Its current image file is the first placeholder extracted from DOCX V4.0.
  \item \textbf{Fig. 2, ``Three-tier deployment map for quantum-safe session establishment and passkey authentication,''} carries the node, gateway, and cloud/server placement content discussed in Section III-D. Its current image file is the second placeholder extracted from DOCX V4.0.
\end{itemize}

The figure numbering, captions, and surrounding references have been reduced to exactly two figures, with no obsolete manuscript references to Figs. 3 or 4. This completed structural revision allows the final artwork to replace the two placeholder files without changing the paper's argument or section organization.

\section*{Response to technical quality and soundness}

\textcolor{reviewblue}{\textbf{Reviewer assessment:}} \emph{Fair (Some technical issues or missing baselines, but potentially fixable).}

\textbf{Response.} The review did not identify a specific missing baseline, and no additional physical-board or accelerator experiment was available for this revision. We therefore did not introduce an unmeasured comparison. Instead, the revised manuscript makes the evaluation boundary explicit in Section IV: the reported values are software-centric QEMU prototype measurements used as placement evidence, not cycle-accurate hardware characterization. Section IV-C also states that Linux/QEMU resident-set size is not an ESP32 SRAM-footprint measurement. The limitations paragraph identifies physical-board power, radio throughput, flash wear, timing leakage, and comparisons with existing PQC IoT architectures as future work.

Within that stated scope, the technical presentation has been tightened in three places. Table I connects each primitive to a deployment tier and rationale. Table II reports the measured elapsed time, resident memory, payload size, and recommended tier without changing the underlying values. Section IV-D links the observations to buffer sizing, gateway admission control, and candidate acceleration kernels. These revisions clarify what the evidence supports and avoid extending the conclusions beyond the prototype.

\section*{Response to originality and novelty}

\textcolor{reviewblue}{\textbf{Reviewer assessment:}} \emph{Fair (Incremental improvement or a minor new application).}

\textbf{Response.} We have not reframed the work as a new cryptographic primitive. The revised Introduction states the contribution more narrowly: (1) a layered node/gateway/server architecture for PQC-enabled IoT, (2) a PQC-aware passkey metadata and hybrid-policy model, and (3) a measurement-driven rule for placing ML-KEM, ML-DSA, and SLH-DSA. Section III and the two revised figures are organized around these system-level contributions. This wording separates the paper's architectural contribution from the standardized algorithms and software implementations on which the prototype relies.

\section*{Summary of manuscript changes}

\footnotesize
\renewcommand{\arraystretch}{1.08}
\begin{tabularx}{\textwidth}{@{}p{0.22\textwidth}Yp{0.24\textwidth}@{}}
\toprule
\textbf{Reviewer concern} & \textbf{Revision} & \textbf{Location} \\
\midrule
AI-generated appearance in submitted Figs. 1 and 4 & Consolidated the four-figure structure to two figures; final author redraw is still pending in this working draft & Placeholder Figs. 1 and 2; Sections III-C--D \\
Low quality of submitted Figs. 2 and 3 & Retained the two DOCX V4.0 images only as replaceable placeholders; no final-quality claim is made & Placeholder Figs. 1 and 2 \\
Technical scope and possible missing baselines & Preserved the measured values, clarified the QEMU/prototype boundary, and identified unperformed board-level comparisons as future work & Section IV-A--D; Conclusion \\
Incremental novelty & Stated the contribution as tiered placement, migration metadata/policy, and measurement-guided system design rather than a new primitive & Introduction; Section III-A; Conclusion \\
\bottomrule
\end{tabularx}
\normalsize

\vspace{0.6em}
We thank the reviewer again. The manuscript text, figure locations, numbering, and captions are ready for the final author-supplied artwork. This response must not be submitted until the two placeholders have been replaced and the figure-quality statements above have been updated to describe the delivered files accurately.

\end{document}
